\documentclass{article}
\usepackage{mathtools} 
\usepackage{fontspec}
\usepackage[UTF8]{ctex}
\usepackage{amsthm}
\usepackage{mdframed}
\usepackage{xcolor}
\usepackage{amssymb}
\usepackage{amsmath}


% 定义新的带灰色背景的说明环境 zremark
\newmdtheoremenv[
  backgroundcolor=gray!10,
  % 边框与背景一致，边框线会消失
  linecolor=gray!10
]{zremark}{说明}

% 通用矩阵命令: \flexmatrix{矩阵名}{元素符号}{行数}{列数}
\newcommand{\flexmatrix}[4]{
  \[
  #1 = \begin{pmatrix}
    #2_{11}     & #2_{12}     & \cdots & #2_{1#4}   \\
    #2_{21}     & #2_{22}     & \cdots & #2_{2#4}   \\
    \vdots      & \vdots      & \ddots & \vdots     \\
    #2_{#31}    & #2_{#32}    & \cdots & #2_{#3#4}
  \end{pmatrix}
  \]
}

% 简化版命令（默认矩阵名为A，元素符号为a）: \quickmatrix{行数}{列数}
\newcommand{\quickmatrix}[2]{\flexmatrix{A}{a}{#1}{#2}}

\begin{document}
\title{习题 5.1}
\author{张志聪}
\maketitle

\section*{4}

易得
\begin{align*}
  f'_+(3) = 6
\end{align*}
左导数为
\begin{align*}
  \lim\limits_{x \to 3^-} \frac{ax + b - f(3)}{x - 3}
   & = \lim\limits_{x \to 3^-} \frac{ax + b - 9}{x - 3}            \\
   & = \lim\limits_{x \to 3^-} \frac{a(x - 3) + 3a + b - 9}{x - 3} \\
   & = \lim\limits_{x \to 3^-} a + \frac{3a + b - 9}{x - 3}        \\
\end{align*}
因为$f$在$x=3$处可导，所以左导数是存在的，所以
\begin{equation*}
  \begin{cases}
    a = 6 \\
    3a + b - 9 = 0
  \end{cases}
\end{equation*}
所以
\begin{equation*}
  \begin{cases*}
    a = 6 \\
    b = -9
  \end{cases*}
\end{equation*}

\section*{9}

\begin{itemize}
  \item (1)

        \begin{align*}
          f'(x) & = cos x + sin x                                                  \\
                & = \sqrt{2} (\frac{\sqrt{2}}{2} cos x + \frac{\sqrt{2}}{2} sin x) \\
                & = \sqrt{2} sin (\frac{\pi}{4} + x)
        \end{align*}
        所以，$x \in \{\frac{3\pi}{4} + k\pi | k \in \mathbb{N} \}, f'(x) = 0$。
\end{itemize}

\section*{12}

由题设，我们有
\begin{align*}
  f(0) = f(0 + 0) = f(0) f(0)
\end{align*}
于是可得$f(x) = 0$或$f(x) = 1$。

如果$f(0) = 0$，那么，$\forall x \in \mathbb{R}$，我们有
\begin{align*}
  f(x + 0) = f(x) f(0) = 0
\end{align*}
即$f$是常值函数，任意位置的导数都是0，与题设$f'(0) = 1$矛盾，
因此，$f(0) = 1$。

$\forall x_0 \in \mathbb{R}$，导数为
\begin{align*}
  \lim\limits_{x \to x_0} \frac{f(x) - f(x_0)}{x - x_0}
  & = \lim\limits_{x \to x_0} \frac{f(x - x_0 + x_0) - f(x_0)}{x - x_0}                              \\
  & = \lim\limits_{x \to x_0} \frac{f(x - x_0)f(x_0) - f(x_0)}{x - x_0}                              \\
  & = \lim\limits_{x \to x_0} f(x_0) \frac{f(x - x_0) - 1}{x - x_0}                                  \\
  & = \lim\limits_{x \to x_0} f(x_0) \frac{f(x - x_0) - f(0)}{(x - x_0) - 0}                         \\
  & = \lim\limits_{x \to x_0} f(x_0) \lim\limits_{x \to x_0} \frac{f(x - x_0) - f(0)}{(x - x_0) - 0} \\
  & = f(x_0) \cdot 1 \\
  & = f(x_0)
\end{align*}
由$x_0$的任意性可知，$f'(x) = f(x)$
\end{document}